\chapter{Родитель композитного
\texorpdfstring{$SU(2)$}{SU(2)}-синглетного полюса}
\label{ch:v10-discrete-resolution-su2-composite-singlet-pole}

Два активных фундаментальных дублета дают разложение
\begin{equation}
 \mathbf2\otimes\mathbf2=\mathbf1\oplus\mathbf3.
\end{equation}
Нормированный синглет
\begin{equation}
 |S\rangle=\frac{|\uparrow\downarrow\rangle-
 |\downarrow\uparrow\rangle}{\sqrt2}
\end{equation}
задаёт проектор $P_S=|S\rangle\langle S|$ ранга один. Он совпадает с
антисимметризатором $(I-\mathsf S)/2$, аннулируется полными генераторами
$J_a=T_a\otimes I+I\otimes T_a$ и поэтому является точным активным
gauge-singlet проектором.

Полный Casimir удовлетворяет
\begin{equation}
 \sum_aJ_a^2=2(I-P_S).
\end{equation}
Следовательно, условный родитель связывания
\begin{equation}
 H_{\rm bind}=I-P_S
\end{equation}
имеет спектр $0^{(1)},1^{(3)}$ и единственный синглетный ground state.
Тем самым ранговый барьер одночастичного K43-сектора снят на композитном
двухдублетном носителе.

Но представление не выбирает физический bound state. Шестнадцать копий
дублета дают $16^2=256$ gauge-singlet flavor-пар; конкретная flavor-линия
не выбрана. Унаследованный двухчастичный родитель равен нулю на gauge-паре
и не создаёт singlet--triplet splitting. Запись условного пропагатора
\begin{equation}
 \Gamma_S^{(2)}(q)=q-e^{-64\pi^2/3}
\end{equation}
даёт простой полюс с единичным вычетом, но его масса пока внесена как цель.

\begin{theorem}[Активный композитный синглет существует]
Антисимметричная компонента двух AF-дублетов является каноническим
$SU(2)$-инвариантным проектором ранга один. Условный положительный родитель
может выбрать её единственным ground state.
\end{theorem}

\begin{nogocriterion}[Кинематика синглета не является связыванием]
Разложение $2\otimes2=1\oplus3$ не выводит привлекательное ядро, flavor-
линию или массовый коэффициент. Без ненулевого singlet--triplet splitting
и полюса четырёхточечной функции композитная архитектура остаётся условной.
\end{nogocriterion}

\begin{protocol}\begin{enumerate}
\item условная архитектура: \textbf{$12/12$};
\item точный gauge-singlet проектор: \textbf{$6/6$};
\item унаследованное связывание: \textbf{$0/2$};
\item физическое происхождение: \textbf{$0/3$};
\item следующий гейт: \textbf{аудит связывающих ядер}.
\end{enumerate}\end{protocol}

Машинный сертификат сохранён в
\path{s2t_v10_cell_birth_four_volume_nonequilibrium_bath_discrete_resolution_transmuted_mode_asymptotically_free_su2_gauge_singlet_composite_pole_parent_origin_gate_results.json}.